% EM structure and teaching revision 20260928
\subsection{电像构造的验收与能量系数}
电像法的替代电荷只能放在不求解的区域。构造完成后依次检查：真实电荷位置的奇性、导体表面的等势条件、远处或其他边界条件。满足这些条件后，唯一性定理才保证指定区域内的电势正确；替代区域内部的场没有对应的物理意义。
点电荷 $q$ 距无限接地平面 $a$，像电荷为平面另一侧的 $-q$。真实电荷所受吸力大小为 $kq^2/(4a^2)$，缓慢从远处移到 $a$ 时体系能量变化为
\[
 W(a)=\int_{\infty}^{a}\frac{kq^2}{4s^2}\dd s
 =-\frac{kq^2}{4a}.
\]
这里外力沿离开平面的方向，故外力沿坐标 $s$ 的分量为正。若把真电荷和像电荷误当成两个独立的真实点电荷，会写出 $-kq^2/(2a)$，多出一倍。
两种处理的区别在于真实导体表面的感应电荷会随 $q$ 的移动重新分布。在线性接地边界下，组装给出 $W=\tfrac12q\varphi_{\mathrm{ind}}$；有另加外场或固定偏置时，外加场能与感应能应分别计账，不能把所有电势项都一律乘以二分之一。

